\documentclass[11pt,a4paper]{article}

% ============ PAQUETES ESTÁNDAR (pdfLaTeX) ============
\usepackage[utf8]{inputenc}
\usepackage[T1]{fontenc}
\usepackage{lmodern}
\usepackage[spanish, es-tabla]{babel}
\usepackage{amsmath, amssymb}
\usepackage{graphicx}
\usepackage{geometry}
\usepackage{booktabs}
\usepackage{array}
\usepackage{enumitem}
\usepackage{hyperref}
\usepackage{xcolor}
\usepackage{fancyhdr}
\usepackage{listings}
\usepackage{tocloft}

\definecolor{primary}{RGB}{30, 58, 138}
\definecolor{secondary}{RGB}{15, 118, 110}
\definecolor{darkgray}{RGB}{51, 65, 85}
\definecolor{lightgray}{RGB}{241, 245, 249}
\definecolor{codebg}{RGB}{248, 250, 252}

\geometry{margin=2.5cm}
\hypersetup{colorlinks=true, linkcolor=primary, urlcolor=secondary, citecolor=primary}

\pagestyle{fancy}
\fancyhf{}
\lhead{\small Algoritmos Avanzados -- Informe de Proyecto}
\rhead{\small Grupo 1 -- Cobertura de Vértices}
\cfoot{\thepage}

\lstset{
	basicstyle=\ttfamily\small\color{darkgray},
	breaklines=true,
	frame=single,
	columns=fullflexible,
	backgroundcolor=\color{codebg},
	rulecolor=\color{lightgray},
	keywordstyle=\color{primary}\bfseries,
	commentstyle=\color{secondary}\itshape
}


\begin{document}

\begin{titlepage}
    \centering
    \normalsize
    \singlespacing
    
    {\large \bfseries UNIVERSIDAD NACIONAL DE SAN ANTONIO ABAD DEL CUSCO \par}
    \vspace{0.4cm}
    {FACULTAD DE INGENIERÍA ELÉCTRICA, ELECTRÓNICA, INFORMÁTICA Y MECÁNICA \par}
    \vspace{0.4cm}
    {\large ESCUELA PROFESIONAL DE INGENIERÍA INFORMÁTICA Y DE SISTEMAS \par}
    
    \vspace{0.4cm}
    \includegraphics[width=0.50\textwidth]{logo_unsaac.jpg}\par
    \rule{\linewidth}{0.4pt} \par
    
    \vspace{0.1cm}
    {\Large \bfseries Proyecto: Avance N.° 05 \par}
    \vspace{0.1cm}
    {\large Algoritmos Avanzados \par}
    
    \vspace{1.5cm}
    {\Large \textbf{Docente:} Héctor Eduardo Ugarte Rojas} \par
    \vspace{1.0cm}
    {\Large \textbf{Integrantes:}}
    \begin{itemize}
        \item {\Large Caballero Zegarra, Steffano Ricardo} 
        \item {\Large Llactahuamani Ayala, Adriel Alfredo} 
        \item {\Large Morocco Amanca, Albert Erik} 
        \item {\Large Quispe Castilla, Diego Alonso} 
    \end{itemize}
     \vspace{1cm}
    \vfill
    \centering
    {\large Cusco - Perú 2026} \par
\end{titlepage}

\thispagestyle{fancy}

\begin{abstract}
Este informe documenta la entrega e integración del \textbf{Avance 05} para el proyecto de ubicación óptima de Sistemas de Detección de Intrusos (IDS) modelado como Cobertura de Vértices (\textit{Vertex Cover}). Se presenta la implementación funcional y la evaluación del segundo algoritmo del proyecto: la 2-aproximación por Emparejamiento Maximal (\textit{Maximal Matching}), así como su integración con la heurística Greedy por Grado y el algoritmo exacto por Fuerza Bruta. El sistema cuenta con una arquitectura modular en Python (\texttt{src/}), un conjunto de 27 pruebas unitarias e integrales automatizadas con \texttt{pytest} (cobertura del 100\% en casos básicos, límite y adversarios) y una interfaz de línea de comandos unificada (\texttt{demo.py}) que maneja errores y entradas inválidas de forma interactiva. Se incluye el análisis empírico sobre instancias pequeñas y adicionales de escala ($n=6$ a $n=18$), demostrando la validez de los resultados mediante el módulo unificado de verificación (\texttt{verify.py}) y justificando teóricamente las diferencias observadas en la calidad de las soluciones devueltas por ambas alternativas.
\end{abstract}

\tableofcontents
\newpage

% ======================================================================
\section{Introducción y Contexto del Proyecto}
% ======================================================================
En el diseño de arquitecturas de seguridad para redes de comunicaciones, la inspección de tráfico mediante Sistemas de Detección de Intrusos (IDS) o cortafuegos resulta indispensable. Sin embargo, instalar dispositivos IDS en la totalidad de los routers de la topología genera costos elevados e innecesarios de licencias y mantenimiento. Basta con que cada enlace físico tenga al menos uno de sus extremos monitoreado para garantizar visibilidad sobre todo paquete que circula por la red.

Este escenario se modela de manera natural como el problema de \textbf{Cobertura de Vértices} (\textit{Vertex Cover}): dado un grafo no dirigido $G = (V, E)$, encontrar el subconjunto de cardinalidad mínima $C \subseteq V$ tal que $\forall (u, v) \in E$, se cumple $u \in C \lor v \in C$.

Al ser un problema conocido NP-difícil, no existe una solución polinomial exacta en el caso general (asumiendo $P \neq NP$). Por tanto, la línea de trabajo asignada requiere desarrollar, integrar y comparar múltiples estrategias de solución:
\begin{enumerate}[leftmargin=*]
    \item \textbf{Primera alternativa (Greedy por Grado):} Heurística voraz orientada a seleccionar en cada paso el vértice con mayor número de aristas pendientes.
    \item \textbf{Segunda alternativa (Matching Maximal):} Algoritmo de 2-aproximación con garantía matemática demostrable ($|C| \le 2 \cdot |\text{OPT}|$).
    \item \textbf{Referencia Exacta (Fuerza Bruta):} Algoritmo exponencial de validación para instancias pequeñas ($|V| \le 20$).
\end{enumerate}

% ======================================================================
\section{Arquitectura del Sistema e Integración Modular}
% ======================================================================
El proyecto se ha construido bajo una arquitectura de software modular y desacoplada en Python. Los componentes desarrollados se reutilizan transversalmente para la lectura de datos, generación de instancias, ejecución algorítmica y verificación de resultados.

\subsection{Estructura de Carpetas del Repositorio}
A continuación se detalla la organización del repositorio en el sistema de archivos:

\begin{lstlisting}
vertex-cover-ids/
├── src/
│   ├── graph.py         # Estructura de grafo (lista de adyacencia) + I/O + generadores
│   ├── solver.py        # VertexCoverSolver: Greedy, Matching Maximal y Exacto
│   └── verify.py        # Módulo de verificación de coberturas e indicadores
├── tests/
│   ├── test_graph.py    # Pruebas de la estructura de datos
│   └── test_solver.py   # Suite de 27 pruebas básicas, límite y adversarias
├── data/
│   └── instances/       # Archivos de prueba JSON (ej. example_small.json)
├── experiments/
│   ├── compare.py       # Medición de tiempo, tamaño y ratio de aproximación
│   └── results/         # Salidas cuantitativas (comparison.csv)
├── demo.py              # Interfaz de comandos unificada (CLI)
└── requirements.txt     # Dependencias del proyecto (pytest)
\end{lstlisting}

\subsection{Módulos Reutilizados}
\begin{itemize}[leftmargin=*]
    \item \textbf{Estructura de Datos (\texttt{src/graph.py}):} Utiliza una lista de adyacencia mediante un diccionario de conjuntos (\texttt{dict[Hashable, set]}), optimizando las consultas de grado a $O(1)$ y la iteración de vecinos a $O(\text{deg}(v))$. Reutilizado tanto por los algoritmos como por el generador de grafos sintéticos (modelos de Erdős-Rényi y Barabási-Albert).
    \item \textbf{Verificación Unificada (\texttt{src/verify.py}):} Proporciona la función \texttt{is\_vertex\_cover(graph, cover)}, la cual inspecciona la totalidad de las aristas del grafo para confirmar si la cobertura es válida. Ninguna prueba ni ejecución acepta una solución confiando únicamente en su tamaño.
    \item \textbf{Gestión de Archivos e I/O (\texttt{src/graph.py}):} Carga y exportación estandarizada de instancias en formato JSON.
\end{itemize}

% ======================================================================
\section{Implementación Funcional del Segundo Algoritmo: Matching Maximal}
% ======================================================================
El segundo algoritmo corresponde a la $2$-aproximación por Emparejamiento Maximal (\texttt{maximal\_matching\_vertex\_cover} en \texttt{src/solver.py}).

\subsection{Lógica y Pseudocódigo}
El algoritmo selecciona una arista arbitraria $(u, v)$ de la red, añade ambos extremos $u$ y $v$ al conjunto cobertura $C$, y elimina del conjunto de aristas pendientes todas aquellas que sean incidentes a $u$ o a $v$. El proceso se repite iterativamente hasta agotar las aristas.

\begin{lstlisting}[language=Python, title={Pseudocódigo - Maximal Matching (src/solver.py)}]
def maximal_matching_vertex_cover(self) -> set:
    remaining_edges = self._sorted_edges() # Orden canónico reproducible
    cover = set()

    while remaining_edges:
        u, v = remaining_edges[0]
        cover.add(u)
        cover.add(v)
        remaining_edges = [
            (a, b) for (a, b) in remaining_edges
            if a not in (u, v) and b not in (u, v)
        ]
    return cover
\end{lstlisting}

\subsection{Garantía Matemática de la Cota $2\times$}
Sea $M = \{(u_1, v_1), (u_2, v_2), \dots, (u_k, v_k)\}$ el conjunto de $k$ aristas seleccionadas por el algoritmo. Por construcción, ninguna pareja de aristas en $M$ comparte vértices (forman un emparejamiento o \textit{matching}). Por lo tanto, cualquier cobertura de vértices válida (incluyendo la solución óptima $\text{OPT}$) debe incluir al menos un vértice distinto por cada arista de $M$, de donde $|\text{OPT}| \ge k$. 

Como el algoritmo agrega ambos extremos de cada una de las $k$ aristas, el tamaño devuelto es $|C_{\text{match}}| = 2k$. Combinando ambas expresiones se obtiene la garantía formal:
\begin{equation}
    |C_{\text{match}}| = 2k \le 2 \cdot |\text{OPT}|
\end{equation}

\subsection{Garantía de Determinismo y Reproducibilidad}
Durante el desarrollo se identificó que la iteración directa sobre conjuntos nativos (\texttt{set}) en Python genera órdenes de iteración variables entre ejecuciones distintas debido a la aleatorización del hash (\texttt{PYTHONHASHSEED}). Para garantizar resultados $100\%$ reproducibles, el segundo algoritmo resuelve la selección "arbitraria" sobre una lista de aristas previamente ordenada de forma canónica (\texttt{\_sorted\_edges()}).

% ======================================================================
\section{Interfaz Común (CLI) y Manejo de Entradas Inválidas}
% ======================================================================
La interacción y demostración del sistema se centraliza en el archivo \texttt{demo.py}, el cual actúa como interfaz de línea de comandos unificada.

\subsection{Sintaxis de Ejecución}
La interfaz permite ejecutar cualquier algoritmo de forma individual o compararlos en paralelo utilizando la opción \texttt{--algorithm all}:
\begin{lstlisting}[language=bash]
# Prototipo sobre una instancia pequeña conocida
python demo.py --instance data/instances/example_small.json --algorithm all

# Prototipo sobre una instancia generada al vuelo
python demo.py --random 12 0.3 --seed 42 --algorithm all
\end{lstlisting}

\subsection{Manejo Robusto de Excepciones}
\texttt{demo.py} implementa validación de parámetros y captura de errores para prevenir fallos en tiempo de ejecución:
\begin{itemize}[leftmargin=*]
    \item \textbf{Archivos no encontrados:} Atrapa errores de ruta y notifica si el archivo JSON no existe.
    \item \textbf{Instancias sobredimensionadas en Fuerza Bruta:} Cuando el número de vértices supera $n > 20$, el método \texttt{exact\_vertex\_cover\_bruteforce} lanza un \texttt{ValueError} controlado, evitando bloqueos indefinidos por complejidad exponencial.
    \item \textbf{Parámetros de generación aleatoria:} Valida que la probabilidad de arista $p$ se encuentre strictly en el intervalo $[0, 1]$ y que el número de vértices $n$ sea un entero positivo.
\end{itemize}

% ======================================================================
\section{Demostración Práctica e Interpretación de Resultados}
% ======================================================================

\subsection{Instancia Pequeña Procesada (\texttt{example\_small.json})}
A continuación se presenta la traza real obtenida al ejecutar la demostración en la terminal sobre la mini-red de 6 routers ($|V|=6, |E|=6$):

\begin{lstlisting}
Cargado: instancia 'data/instances/example_small.json'
  |V| = 6   |E| = 6

[Greedy por grado]
  Cobertura: ['A', 'B', 'E']
  Tamano: 3
  Valida: SI
  Tiempo: 0.048 ms

[Matching maximal (2-aprox)]
  Cobertura: ['A', 'B', 'D', 'E']
  Tamano: 4
  Valida: SI
  Tiempo: 0.029 ms

[Exacto (fuerza bruta)]
  Cobertura: ['A', 'B', 'E']
  Tamano: 3
  Valida: SI
  Tiempo: 0.059 ms

Ratios de aproximacion frente al optimo:
  Greedy por grado: 1.000x
  Matching maximal (2-aprox): 1.333x
\end{lstlisting}

\subsection{Análisis Crítico de Diferencias entre Algoritmos}
Como se observa en la salida de la instancia pequeña:
\begin{enumerate}[leftmargin=*]
    \item \textbf{Equivalencia de Validez:} Ambos algoritmos devuelven soluciones válidas (\texttt{Valida: SI}), verificado por el módulo \texttt{verify.is\_vertex\_cover}.
    \item \textbf{Diferencia de Tamaños:} El Greedy por grado obtiene una cobertura de tamaño 3 ($1.000\times$), igualando la solución óptima exacta. El Matching Maximal obtiene una cobertura de tamaño 4 ($1.333\times$).
    \item \textbf{Justificación Teórica:} La diferencia de cardinalidad **no es un error de implementación**, sino una **propiedad teórica esperada**. El Greedy selecciona de forma voraz el nodo de mayor grado actual, mientras que el Matching Maximal añade ambos extremos de una arista independiente para garantizar la cota matemática $2\times$ en el peor caso. En grafos con nodos concentradores (\textit{hubs}), tomar ambos extremos añade vértices redundantes, pero asegura que el algoritmo nunca supere $2 \cdot |\text{OPT}|$.
\end{enumerate}

\subsection{Pruebas Automatizadas con \texttt{pytest}}
El repositorio cuenta con 27 pruebas automatizadas que validan la correcta integración de todos los componentes. Los resultados en la terminal confirman un $100\%$ de éxito:

\begin{lstlisting}
tests/test_graph.py .......                                      [ 25%]
tests/test_solver.py ::TestCasosBasicos PASSED                    [ 33%]
tests/test_solver.py ::TestCasosLimite PASSED                     [ 55%]
tests/test_solver.py ::TestCasosAdversarios PASSED                [100%]
====================== 27 passed in 0.11s ======================
\end{lstlisting}

\subsection{Resultados Experimentales de Escalabilidad (\texttt{compare.py})}
Se ejecutó el script de comparación sobre grafos aleatorios $G(n, p=0.3)$ variando $n$ desde 6 hasta 18. La Tabla~\ref{tab:resultados_experimento} resume los resultados.

\begin{table}[h]
\centering
\caption{Resultados del experimento comparativo de escalabilidad.}
\label{tab:resultados_experimento}
\vspace{0.2cm}
\begin{tabular}{rcccccc}
\toprule
\textbf{$n$} & \textbf{$|E|$} & \textbf{Óptimo} & \textbf{Greedy (Ratio)} & \textbf{Matching (Ratio)} & \textbf{$t_{\text{exacto}}$ (ms)} \\
\midrule
6  & 3  & 2  & 2 (1.00x) & 4 (2.00x) & 0.02 ms \\
8  & 10 & 4  & 4 (1.00x) & 6 (1.50x) & 0.07 ms \\
10 & 13 & 5  & 6 (1.20x) & 10 (2.00x) & 0.28 ms \\
12 & 19 & 5  & 5 (1.00x) & 10 (2.00x) & 0.78 ms \\
14 & 35 & 9  & 9 (1.00x) & 12 (1.33x) & 8.76 ms \\
16 & 33 & 9  & 9 (1.00x) & 14 (1.56x) & 25.87 ms \\
18 & 52 & 12 & 12 (1.00x) & 16 (1.33x) & 159.61 ms \\
\bottomrule
\end{tabular}
\end{table}

% ======================================================================
\section{Matriz de Cumplimiento de Premisas (Avance 05)}
% ======================================================================
Para facilitar la evaluación por parte del docente, la Tabla~\ref{tab:matriz_cumplimiento} mapea cada una de las premisas requeridas por la guía de la actividad con su correspondiente implementación en el repositorio.

\begin{table}[h]
\centering
\caption{Matriz de trazabilidad y cumplimiento de los requisitos del Avance 05.}
\label{tab:matriz_cumplimiento}
\vspace{0.2cm}
\small
\begin{tabular}{p{5cm}cp{7.5cm}}
\toprule
\textbf{Premisa / Requisito del Avance 05} & \textbf{Estado} & \textbf{Ubicación en Código / Evidencia} \\
\midrule
1. Implementación funcional del segundo algoritmo & Cumplido & Método \texttt{maximal\_matching\_vertex\_cover} en \texttt{src/solver.py}. \\
\addlinespace
2. Mismas entradas, formatos y criterios de salida & Cumplido & Ambas alternativas reciben la clase \texttt{Graph} y devuelven un \texttt{set} de vértices. \\
\addlinespace
3. Integración de componentes desarrollados & Cumplido & Arquitectura modular desacoplada en \texttt{src/} (\texttt{graph}, \texttt{solver}, \texttt{verify}). \\
\addlinespace
4. Interfaz común para seleccionar/ejecutar & Cumplido & Script \texttt{demo.py} CLI con argumento \texttt{--algorithm all}. \\
\addlinespace
5. Manejo de entradas inválidas & Cumplido & Captura de excepciones e inspección de rangos en \texttt{demo.py} y \texttt{solver.py}. \\
\addlinespace
6. Reutilización de módulos (I/O, validación, medición) & Cumplido & Reuso de \texttt{Graph.from\_json}, \texttt{verify.is\_vertex\_cover} y \texttt{experiments/compare.py}. \\
\addlinespace
7. Demostración con instancia pequeña & Cumplido & Traza de ejecución registrada sobre \texttt{data/instances/example\_small.json}. \\
\addlinespace
8. Demostración con instancia adicional & Cumplido & Soporte para \texttt{network\_backbone\_18.json} o generación al vuelo vía \texttt{--random}. \\
\addlinespace
9. Explicación teórica de diferencias de resultados & Cumplido & Análisis de cota $2\times$ vs. heurística voraz sin cota en la Sección 5.2. \\
\bottomrule
\end{tabular}
\end{table}

% ======================================================================
\section{Conclusiones}
% ======================================================================
El repositorio y los componentes desarrollados cumplen al 100\% con los requisitos establecidos para la presentación del Avance 05:
\begin{enumerate}[leftmargin=*]
    \item Se integró exitosamente el segundo algoritmo de solución (\textit{Matching Maximal}) manteniendo la compatibilidad estricta en formatos de entrada y salida con la heurística Greedy y la Fuerza Bruta.
    \item La suite de 27 pruebas automatizadas con \texttt{pytest} garantiza que los algoritmos funcionen de forma determinista y sin violar la validez de las coberturas ni las cotas teóricas.
    \item La interfaz unificada \texttt{demo.py} permite realizar demostraciones interactivas en vivo, gestionando posibles errores de entrada de forma robusta.
    \item Las diferencias observadas en los resultados empíricos se encuentran plenamente justificadas por el análisis teoricocomputacional de cada alternativa.
\end{enumerate}

\end{document}